\documentclass[11pt,a4paper]{article}
\usepackage[UTF8]{ctex}
\usepackage{amsmath,amssymb,bm}
\usepackage{geometry}
\usepackage{hyperref}
\geometry{margin=2.2cm}

\title{iphase 时差公式（与当前代码一致）}
\author{}
\date{\today}

\begin{document}
\maketitle

\section{符号约定}
\begin{itemize}
\item $\Delta t$：时差
\item $V_p, V_s$：末段介质中的 P/S 波速度
\item $h$：末段垂向厚度（或等效厚度）
\item $L$：末段路径长度
\item $p=\mathrm{d}t/\mathrm{d}x$：由走时曲线斜率估计的慢度
\item $x_{\mathrm{conv}}$：转换点位置，$x_{\mathrm{rec}}$：接收点位置
\end{itemize}

\section{观测时差定义}
\subsection{PPS-PPP}
严格配对模式：
\begin{equation}
\Delta t_{\mathrm{PPS-PPP}} = t_{\mathrm{PPS}} - t_{\mathrm{PPP}}.
\end{equation}

拟合取值模式（默认）：
\begin{equation}
\Delta t_{\mathrm{PPS-PPP}} = t_{\mathrm{PPS,fit}}(x) - t_{\mathrm{PPP,fit}}(x).
\end{equation}

\subsection{PSS-PSP}
严格配对模式：
\begin{equation}
\Delta t_{\mathrm{PSS-PSP}} = t_{\mathrm{PSS}} - t_{\mathrm{PSP}}.
\end{equation}

拟合取值模式（默认）：
\begin{equation}
\Delta t_{\mathrm{PSS-PSP}} = t_{\mathrm{PSS}} - t_{\mathrm{PSP,fit}}(x).
\end{equation}

\section{理论核心式}
两类时差在代码中的统一核心为
\begin{equation}
\Delta t = L\left(\frac{1}{V_s} - \frac{1}{V_p}\right).
\end{equation}
在通常条件 $V_s<V_p$ 下，$\Delta t$ 为正。

\section{PPS-PPP 主理论（1D）}
由 PPP 拟合走时曲线估计
\begin{equation}
p=\frac{\mathrm{d}t}{\mathrm{d}x}, \qquad
\sin\theta_p = |p|V_p.
\end{equation}
路径长度取
\begin{equation}
L_{ppp}=\frac{h}{\cos\theta_p}
=\frac{h}{\sqrt{1-(|p|V_p)^2}}.
\end{equation}
代入得
\begin{equation}
\Delta t_{\mathrm{PPS-PPP}}^{\mathrm{main}}
=L_{ppp}\left(\frac{1}{V_s}-\frac{1}{V_p}\right).
\end{equation}

\section{1D 诊断图4的 PPS-PPP 三条理论线}
当前代码在图4中绘制三条口径：
\begin{align}
\Delta t_{\mathrm{main}} &= \frac{L_{pps}}{V_s} - \frac{L_{ppp}}{V_p}, \\
\Delta t_{\mathrm{alt1}} &= \frac{L_{ppp}}{V_s} - \frac{L_{ppp}}{V_p}, \\
\Delta t_{\mathrm{alt2}} &= \frac{L_{pps}}{V_s} - \frac{L_{pps}}{V_p}.
\end{align}
其中主线为第一个式子；其余两条为对照线。

\section{1D 诊断图4的 PSS-PSP 三条理论线}
当前代码同样绘制三条口径：
\begin{align}
\Delta t_{\mathrm{main}} &= \frac{L_{pss}}{V_s} - \frac{L_{psp}}{V_p}, \\
\Delta t_{\mathrm{alt1}} &= \frac{L_{psp}}{V_s} - \frac{L_{psp}}{V_p}, \\
\Delta t_{\mathrm{alt2}} &= \frac{L_{pss}}{V_s} - \frac{L_{pss}}{V_p}.
\end{align}

\section{PSS-PSP 主算法口径（与诊断图不同）}
主算法中，PSS-PSP 使用迭代几何长度：
\begin{equation}
L=\sqrt{h^2+\left(x_{\mathrm{conv}}-x_{\mathrm{rec}}\right)^2},
\end{equation}
再代入
\begin{equation}
\Delta t_{\mathrm{PSS-PSP}} =
L\left(\frac{1}{V_s}-\frac{1}{V_p}\right).
\end{equation}
因此，诊断图中的绿色虚线/点划线只用于形态对照，不等同于主算法反演口径。

\section{结论}
\begin{itemize}
\item 观测定义统一为“后相位减前相位”。
\item 理论核心统一为 $\Delta t=L(1/V_s-1/V_p)$。
\item 不同模块之间差异主要在 $L$ 的构造方式。
\end{itemize}

\end{document}
